% topic: programming
\question บริษัทแห่งหนึ่งได้เขียนโค้ดภาษา Python เพื่อคำนวณภาษีเงินได้ประจำปีของพนักงานตามเกณฑ์ขั้นบันได ดังนี้:
\begin{itemize}
    \item รายได้ น้อยกว่า 30,000 บาท เสียภาษี 0\%
    \item รายได้ ตั้งแต่ 30,000 ถึง 100,000 บาท เสียภาษี 10\%
    \item รายได้ มากกว่า 100,000 บาท เสียภาษี 20\%
\end{itemize}
นักพัฒนาได้ส่งโค้ดนี้เข้าสู่ระบบ:
\begin{verbatim}
if income < 30000:
    return 0
elif income > 30000 and income < 100000:
    return income * 0.10
elif income > 100000:
    return income * 0.20
else:
    return 0
\end{verbatim}
ระบบนี้มีบั๊ก ทำให้พนักงานบางคนไม่ต้องเสียภาษีเลยสักบาท ทั้งที่จริง ๆ แล้วพวกเขาต้องเสียภาษี ค่าที่ทำให้เกิดบั๊กตรงกับข้อใด\newline
\choicesgrid{income = 29999 และ income = 100001}{income = 30000 และ income = 100000}{income = 0 และ income = 50000}{income = -5000 และ income = 100000}
% answer: 2